\subsection{Đề 1}
\caulc
\begin{ex}%[1D1H5-3]
Biểu diễn họ nghiệm của phương trình $\sin2x=1$ trên đường tròn đơn vị ta được bao nhiêu điểm?
\choice
{$1$}
{$8$}
{$4$}
{\True $2$}
\loigiai{
Ta có $\sin2x=1\Leftrightarrow 2x=\dfrac{\pi}{2}+2k\pi\Leftrightarrow x=\dfrac{\pi}{4}+k\pi,\,k\in \mathbb{Z}$.\\
Vậy có $2$ điểm trên đường tròn đơn vị.
}
\end{ex}

\begin{ex}%[1D2H3-5]
Tìm tất cả giá trị của $x$ để ba số $2x-1$, $x$, $2x+1$ theo thứ tự đó lập thành một cấp số nhân.
\choice
{\True $x=\pm \dfrac{1}{\sqrt{3}}$}
{$x=\pm \dfrac{1}{3}$}
{$x=\pm \sqrt{3}$}
{$x=\pm 3$}
\loigiai{
Ba số $2x-1$, $x$, $2x+1$ theo thứ tự đó lập thành một cấp số nhân
\begin{eqnarray*}
&\Rightarrow&x^2=(2x-1)(2x+1)\\
&\Leftrightarrow&x^2=4x^2-1\\
&\Leftrightarrow&3x^2=1\\
&\Leftrightarrow&x=\pm \dfrac{1}{\sqrt{3}}.
\end{eqnarray*}
}
\end{ex}

\begin{ex}%[1D3N1-2]
Giá trị của giới hạn $\lim \dfrac{3n^3-2n+1}{4n^4+2n+1}$ bằng
\choice
{$+\infty$}
{$\dfrac{3}{4}$}
{\True $0$}
{$\dfrac{2}{7}$}
\loigiai{
Ta có
\begin{eqnarray*}
\lim \dfrac{3n^3-2n+1}{4n^4+2n+1}&=&\lim\dfrac{n^3\left(3-\dfrac{2}{n^2}+\dfrac{1}{n^3}\right)}{n^4\left(4+\dfrac{2}{n^3}+\dfrac{1}{n^4}\right)}\\
&=& \lim\left[\dfrac{1}{n}\cdot \dfrac{3-\dfrac{2}{n^2}+\dfrac{1}{n^3}}{4+\dfrac{2}{n^3}+\dfrac{1}{n^4}}\right] \\
&=&0\cdot \dfrac{3-0+0}{4+0+0}=0.
\end{eqnarray*}
}
\end{ex}

\begin{ex}%[1D3N2-2]
	Giá trị của giới hạn $\lim\limits_{x\to 2} (3x^2+7x+11)$ bằng
	\choice
	{\True $37$}
	{$38$}
	{$39$}
	{$40$}
	\loigiai{
Ta có 
$\lim\limits_{x\to 2} (3x^2+7x+11)=3\cdot 2^2+7\cdot 2+11=37$.
}
\end{ex}

\begin{ex}%[1D3N3-3]
Cho hàm số $f(x)=\heva{&x^2+2x+m&\text{khi}\quad x\ge 2\\&3&\text{khi}\quad x<2}$. Giá trị của tham số $m$ để hàm số liên tục tại $x=2$ là
\choice
{$m=3$}
{$m=5$}
{$m=-3$}
{\True $m=-5$}
\loigiai{
Ta có $\heva{&\lim\limits_{x\to 2^{+}}f(x)=f(2)=2^2+2\cdot2+m=8+m\\&\lim\limits_{x\to 2^{-}}f(x)=3.}$\\
Để hàm số liên tục tại $x=2\Rightarrow f(2)=\lim\limits_{x\to 2^{+}}f(x)=\lim\limits_{x\to 2^{-}}f(x)\Leftrightarrow 8+m=3\Leftrightarrow m=-5$.
}
\end{ex}

\begin{ex}%[1D3N3-3]
Số điểm gián đoạn của hàm số $f(x)=\heva{&0{,}5&\text{khi}&\quad x=-1\\&\dfrac{x(x+1)}{x^2-1}&\text{khi}&\quad x\ne -1,x\ne 1\\&1&\text{khi}&\quad x=1}$ là
\choice
{$0$}
{\True $1$}
{$2$}
{$3$}
\loigiai{
Ta có $\heva{&\lim\limits_{x\to 1}f(x)=\lim\limits_{x\to 1}\dfrac{x(x+1)}{x^2-1}=\lim\limits_{x\to 1}\dfrac{x}{x-1}=+\infty\\&f(1)=1.}$\\
Do đó hàm số gián đoạn tại $x=1$.\\
Mặt khác $\heva{&\lim\limits_{x\to -1}f(x)=\lim\limits_{x\to -1}\dfrac{x(x+1)}{x^2-1}=\lim\limits_{x\to -1}\dfrac{x}{x-1}=\dfrac{1}{2}\\&f(-1)=0{,}5.}$\\
Do đó hàm số liên tục tại $x=-1$.\\
Vậy hàm số có $1$ điểm gián đoạn.
}
\end{ex}

\begin{ex}%[1H4N4-2]
Cho hình hộp $ABCD.A'B'C'D'$. Mặt phẳng $(AB'D')$ song song với mặt phẳng nào sau đây?
\choice
{$(BA'C')$}
{\True $(C'BD)$}
{$(BDA')$}
{$(ACD')$}
\loigiai{
\immini{
Ta có $\heva{&B'D'\parallel BD\\&BD\subset (C'BD)}\Rightarrow B'D'\parallel (C'BD)$.\\
Mặt khác $\heva{&AB'\parallel C'D\\&C'D\subset (C'BD)}\Rightarrow AB'\parallel (C'BD)$.\\
Mà $\heva{&B'D',AB'\subset (AB'D')\\&B'D'\cap AB'=B'}$.\\
Vậy $(AB'D')\parallel ((C'BD))$.
}
{\begin{tikzpicture}[scale=1,line join=round,font=\footnotesize ,line cap=round,>=stealth,thick]
	\def\a{3.5}
	\def\h{4.5}
	\path 	(0:0) coordinate (A)
	++(0:\a) coordinate (D)
	++(-130:\a/2) coordinate (C)
	($(A)+(C)-(D)$) coordinate (B)	
	($(A)!0.45!(C)$) coordinate (H)
	($(H)+(90:\h)$) coordinate (A')
	($(A')+(C)-(A)$) coordinate (C')
	($(C')+(B)-(C)$) coordinate (B')
	($(A')+(D)-(A)$) coordinate (D');
	\draw[dashed,thick] 	(B)--(A)--(D)	(A)--(A') (D')--(A)--(B') (B)--(D);
	\draw[thick] (B')--(D')	(C)--(C') 	(D)--(D') 	(B)--(B')  (B)--(C)--(D) (A')--(B')--(C')--(D')--cycle (D)--(C')--(B);
	\foreach \x/\g in {A/180,B/180,C/0,D/0,A'/180,B'/180,C'/0,D'/0}
	\fill[black] 	(\x) circle (1pt)
	($(\g:4mm)+(\x)$) node {$\x$};
%	\draw pic[draw,angle radius=2mm]{right angle=A'--H--C};%Theo chiều dương		
\end{tikzpicture}}
}
\end{ex}

\begin{ex}%[1H4N6-1]
Qua phép chiếu song song, tính chất nào sau đây \textbf{không} được bảo toàn?
\choice
{\True chéo nhau}
{đồng qui}
{song song}
{thẳng hàng}
\loigiai{
Qua phép chiếu song song, tính chất chéo nhau không được bảo toàn.
}
\end{ex}

\begin{ex}%[1D5N1-3]
Điều tra về số tiền mua đồ dùng học tập (đơn vị: nghìn đồng) trong một tháng của $40$ học sinh ta có mẫu số liệu như sau:
\begin{longtable}{|c|c|c|c|c|c|c|c|}
\hline
Số tiền&$[10;15)$&$[15;20)$&$[20;25)$&$[25;30)$&$[30;35)$&$[35;40)$&Cộng\\
\hline
Số học sinh&$2$&$5$&$15$&$8$&$9$&$1$&$N=40$\\
\hline
\end{longtable}
Số trung bình của mẫu số liệu là
\choice
{$22{,}5$}
{\True $25$}
{$25{,}5$}
{$27$}
\loigiai{
Ta có bảng số liệu với giá trị đại diện như sau:
\begin{longtable}{|c|c|c|c|c|c|c|}
	\hline
	Số tiền&$[10;15)$&$[15;20)$&$[20;25)$&$[25;30)$&$[30;35)$&$[35;40)$\\
	\hline
	Giá trị đại diện&$12{,}5$&$17{,}5$&$22{,}5$&$27{,}5$&$32{,}5$&$37{,}5$\\
	\hline
	Số học sinh&$2$&$5$&$15$&$8$&$9$&$1$\\
	\hline
\end{longtable}
Số trung bình $\overline{x}=\dfrac{12{,}5\cdot 2+17{,}5\cdot 5+22{,}5\cdot 15+27{,}5\cdot 8+32{,}5\cdot 9+37{,}5\cdot 1}{40}=25$.
}
\end{ex}

\begin{ex}%[1D5N1-4]
Khảo sát thời gian (đơn vị: phút) tập thể dục của một số học sinh khối $11$ thu được mẫu số liệu ghép nhóm sau:
\begin{longtable}{|c|c|c|c|c|}
\hline
Thời gian (phút)&$[0;20)$&$[20;40)$&$[40;60)$&$[60;80)$\\
\hline
Số học sinh&$5$&$9$&$12$&$10$\\
\hline
\end{longtable}
Nhóm chứa mốt của mẫu số liệu trên là
\choice
{$[0;20)$}
{$[20;40)$}
{\True $[40;60)$}
{$[60;80)$}
\loigiai{
Nhóm chứa mốt là nhóm có tần số lớn nhất.\\
Do đó nhóm chứa mốt là nhóm $[40;60)$ có tần số là $12$.
}
\end{ex}

\begin{ex}%[1D5N2-2]
Một học viện bóng đá điều tra về lứa tuổi của $85$ học viên trẻ đăng kí đầu tiên để tham gia khóa học mới và thu được bảng sau:
\begin{longtable}{|c|c|c|c|c|}
\hline
Nhóm tuổi&$[8;10)$&$[10;12)$&$[12;14)$&$[14;16)$\\
\hline
Số học viên&$14$&$20$&$33$&$18$\\
\hline
\end{longtable}
Nhóm chứa trung vị của mẫu số liệu ghép nhóm trên là
\choice
{$[8;10)$}
{$[10;12)$}
{\True $[12;14)$}
{$[14;16)$}
\loigiai{
Cỡ mẫu $n=85$.\\
Gọi $x_1,x_2,\ldots,x_{85}$ lần lượt là số tuổi của $85$ học viên xếp theo thứ tự không giảm.\\
Do đó trung vị của mẫu số liệu $x_1,x_2,\ldots,x_{85}$ là $M_e=x_{43}\in [12;14)$.
}
\end{ex}

\begin{ex}%[1D5N2-3]
Tổng hợp tiền lương tháng (đơn vị: triệu đồng) của một số nhân văn phòng được ghi lại như sau:
\begin{longtable}{|c|c|c|c|c|}
\hline
Lương tháng (triệu đồng) &$[6;8)$&$[8;10)$&$[10;12)$&$[12;14)$\\
\hline
Số nhân viên&$3$&$6$&$8$&$7$\\
\hline
\end{longtable}
Giá trị của tứ phân vị thứ nhất bằng
\choice
{$11$}
{\True $9$}
{$8$}
{$10$}
\loigiai{
Cỡ mẫu $n=24$.\\
Gọi $x_1,x_2,\ldots,x_{24}$ lần lượt là số tiền lương của $24$ nhân viên xếp theo thứ tự không giảm.\\
Tứ phân vị thứ nhất của mẫu số liệu $x_1,x_2,\ldots,x_{24}$ là $Q_1=\dfrac{1}{2}(x_6+x_7)\in [8;10)$.\\
Do đó giá trị của tứ phân vị thứ nhất là $Q_1=8+\dfrac{\dfrac{24}{4}-3}{6}\cdot (10-8)=9$.
}
\end{ex}

\cauds
\begin{ex}%[1H4N4-1]
Các mệnh đề sau đúng hay sai?
\choiceTF
{\True Hai mặt phẳng phân biệt không cắt nhau thì song song}
{Nếu mặt phẳng này chứa hai đường thẳng cùng song song với mặt phẳng kia thì hai mặt phẳng đó song song với nhau}
{\True Hai mặt phẳng cùng song song với mặt phẳng thứ ba thì song song với nhau}
{Hai mặt phẳng phân biệt cùng song song với một đường thẳng thì song song với nhau}
\loigiai{
\begin{itemchoice}
\itemch{\bf Đúng.}\\
Hai mặt phẳng phân biệt không cắt nhau thì song song.	
\itemch{\bf Sai.}\\
Vì nếu mặt phẳng này chứa hai đường thẳng song song cùng song song với mặt phẳng kia thì hai mặt phẳng đó có thể không song song với nhau.
\itemch{\bf Đúng.}\\
Hai mặt phẳng cùng song song với mặt phẳng thứ ba thì song song với nhau.
\itemch{\bf Sai.}\\
Vì hai mặt phẳng có thể cắt nhau theo giao tuyến song song với đường thẳng đã cho.
\end{itemchoice}
}
\end{ex}

\begin{ex}%[1H4H1-3]
Cho tứ diện $ABCD$. Gọi $I$, $J$ lần lượt là trung điểm của $AD$, $BC$, $M$ là một điểm trên cạnh 
$AB$, $N$ là một điểm trên cạnh $AC$. Khi đó 
\choiceTF
{\True $IJ$ là giao tuyến của hai mặt phẳng $(IBC)$ và $(JAD)$}
{\True $ND$ là giao tuyến của hai mặt phẳng $(MND)$ và $(ADC)$}
{\True $BI$ là giao tuyến của hai mặt phẳng $(BCI)$ và $(ABD)$}
{Giao tuyến của hai mặt phẳng $(IBC)$ và $(DMN)$ song song với đường thẳng $IJ$}
\loigiai{
\begin{center}
\begin{tikzpicture}[scale=1,line join=round,font=\footnotesize ,line cap=round,>=stealth,thick]
	\def\a{4.5}
	\path 	(0:0) coordinate (B)
	++(0:1.2*\a) coordinate (D)
	++(-150:\a) coordinate (C)
	($(B)+(70:\a)$) coordinate (A);
\path ($(A)!0.5!(D)$) coordinate (I);
\path ($(B)!0.5!(C)$) coordinate (J);
\path ($(A)!0.3!(B)$) coordinate (M);
\path ($(A)!0.4!(C)$) coordinate (N);
\coordinate (K) at (intersection of B--I and M--D);
\coordinate (L) at (intersection of C--I and N--D);
	\draw[dashed,thick] 	(B)--(D) (I)--(J) (M)--(D)(B)--(I)(J)--(D);
	\draw[thick] 			(B)--(A)--(D) (A)--(C)
	(B)--(C)--(D) (D)--(N)--(M) (I)--(C) (J)--(A);
	\foreach \x/\g in {A/90,B/180,C/-45,D/0,I/45,J/180,M/180,N/30,K/90,L/-90}
	\fill[black] 	(\x) circle (1.5pt)
	($(\g:3mm)+(\x)$) node {$\x$};
	%Hình chóp S.ABC có SA vuông góc đáy	
\end{tikzpicture}
\end{center}
\begin{itemchoice}
\itemch{\bf Đúng.}\\
Ta có $\heva{&I\in (IBC)\\&I\in AD\\&AD\subset (JAD)}\Rightarrow I\in (IBC)\cap (JAD).\quad (1)$\\
Mặt khác $\heva{&J\in (JAD)\\&J\in BC\\&BC\subset (IBC)}\Rightarrow J\in (IBC)\cap (JAD).\quad (2)$\\
Từ $(1)$ và $(2)\Rightarrow IJ=(IBC)\cap (JAD)$.
\itemch{\bf Đúng.}\\
Ta có $\heva{&ND\subset (MND)\\&ND\subset (ADC)}\Rightarrow ND=(MND)\cap(ADC)$.
\itemch{\bf Đúng.}\\
Ta có $\heva{&BI\subset (BCI)\\&BI\subset (ABD)}\Rightarrow BI=(BCI)\cap(ABD)$.
\itemch{\bf Sai.}\\
Gọi $K=BI\cap MD$ và $L=CI\cap ND\Rightarrow KL=(IBC)\cap (DMN)$.\\
Mà $KL$ không song song $IJ$.\\
Do đó giao tuyến của hai mặt phẳng $(IBC)$ và $(DMN)$ không song song với đường thẳng $IJ$.
\end{itemchoice}
}
\end{ex}

\begin{ex}%[1H4H3-2]
Cho hình chóp $S.ABCD$ có đáy $ABCD$ là hình bình hành. Lấy điểm $M$ trên cạnh $AD$ sao cho 
$AD=3AM$. Gọi $G$, $N$ theo thứ tự là trọng tâm các tam giác $SAB$ và $ABC$. Khi đó
\choiceTF
{Giao tuyến của hai mặt phẳng  $(SAB)$ và $(SCD)$ là đường thẳng đi qua $S$ và song song với $AC$, $BD$}
{$\dfrac{DN}{DB}=\dfrac{1}{3}$}
{\True $MN$ song song với mặt phẳng $(SCD)$}
{\True $NG$ song song với mặt phẳng $(SAC)$}
\loigiai{
\begin{center}
	\begin{tikzpicture}[scale=1,line join=round,font=\footnotesize ,line cap=round,>=stealth,thick]
		\coordinate (A) at (0.5,0);
		\coordinate (B) at (-2,-2);
		\coordinate (D) at (6,0);
		\coordinate (C) at ($(B)+(D)-(A)$);
		\coordinate (S) at ($(A)+(1,5)$);
\path ($(A)!1/3!(D)$) coordinate (M);
\path ($(A)!0.5!(B)$) coordinate (K);
\path ($(A)!0.5!(C)$) coordinate (O);
\path ($(S)!2/3!(K)$) coordinate (G);
\path ($(C)!2/3!(K)$) coordinate (N);
		\draw(S)--(B) (S)--(C) (S)--(D) (B)--(C)--(D);
		\draw[dashed,thin](S)--(A) (A)--(B) (A)--(D)(S)--(K)--(C)--(A) (B)--(D) (G)--(N)--(M);
		\foreach \i/\g in {S/90,A/45,B/-90,C/-90,D/0,M/90,K/-90,G/0,N/-90,O/90}{\draw[fill=black](\i) circle (1.5pt) ($(\i)+(\g:3mm)$) node[scale=1]{$\i$};}
	\end{tikzpicture}
\end{center}
\begin{itemchoice}
	\itemch{\bf Sai.}\\
Giao tuyến của hai mặt phẳng  $(SAB)$ và $(SCD)$ là đường thẳng đi qua $S$ và song song với $AB$, $CD$.
	\itemch{\bf Sai.}\\
	Gọi $O=AC\cap BD$, khi đó $\dfrac{BN}{BO}=\dfrac{2}{3}\Rightarrow \dfrac{BN}{BD}=\dfrac{1}{3}\Rightarrow \dfrac{DN}{DB}=\dfrac{2}{3}$.
	\itemch{\bf Đúng.}\\
Ta có $\heva{&\dfrac{DN}{DB}=\dfrac{2}{3}\quad\text{(cmt)}\\&\dfrac{DM}{DA}=\dfrac{2}{3}\quad\text{(gt)}}\Rightarrow MN\parallel AB\parallel CD$.\\
Mà $CD\subset (SCD)\Rightarrow MN\parallel(SCD)$.
\itemch{\bf Đúng.}\\
Gọi $K$ là trung điểm $AB$.\\
Ta có $\heva{&\dfrac{KG}{KS}=\dfrac{1}{3}\quad(\text{do G là trọng tâm}\,\,\triangle SAB)\\&\dfrac{KN}{KC}=\dfrac{1}{3}\quad(\text{do N là trọng tâm}\,\,\triangle ABC)}\Rightarrow NG\parallel SC$.\\
Mà $SC\subset (SAC)\Rightarrow NG\parallel (SAC)$.
\end{itemchoice}
}
\end{ex}

\begin{ex}%[1D5H2-3]
Khi đo mắt cho một số học sinh khối $11$ ở một trường THPT, nhân viên y tế ghi nhận lại ở bảng sau:
\begin{longtable}{|c|c|c|c|c|c|}
\hline 
Độ cận&$[0{,}25;0{,}75]$&$[0{,}75;1{,}25]$&$[1{,}25;1{,}75]$&$[1{,}75;2{,}25]$&$[2{,}25;2{,}75]$\\
\hline 
Số học sinh&$25$&$32$&$14$&$12$&$7$\\
\hline
\end{longtable}
Khi đó
\choiceTF
{\True Độ cận trung bình của số học sinh trên xấp xỉ $1{,}19$ độ}
{\True $0{,}89$ độ là độ cận mà số học sinh được đo ở trên chiếm nhiều nhất}
{$50\%$ số học sinh được đo có độ cận không vượt quá $0{,}75$ độ}
{\True $50\%$ số học sinh được đo có độ cận từ $0{,}7$ độ đến $1{,}625$ độ}
\loigiai{
Ta có mẫu số liệu ghép nhóm với giá trị đại diện như sau:
\begin{longtable}{|c|c|c|c|c|c|}
	\hline 
	Độ cận&$[0{,}25;0{,}75]$&$[0{,}75;1{,}25]$&$[1{,}25;1{,}75]$&$[1{,}75;2{,}25]$&$[2{,}25;2{,}75]$\\
	\hline 
	Giá trị đại diện&$0{,}5$&$1{,}0$&$1{,}5$&$2{,}0$&$2{,}5$\\
	\hline
	Số học sinh&$25$&$32$&$14$&$12$&$7$\\
	\hline
\end{longtable}
Gọi $x_1,x_2,\ldots,x_{90}$ lần lượt là độ cận của $90$ học sinh được sắp xếp theo thứ tự không giảm.
\begin{itemchoice}
\itemch{\bf Đúng.}\\
Độ cận trung bình $\overline{x}=\dfrac{0{,}5\cdot 25+1\cdot 32+1{,}5\cdot 14+2\cdot 12+2{,}5\cdot 7}{90}\approx 1{,}19$.
\itemch{\bf Đúng}\\
Ta có nhóm chứa mốt là $[0{,}75;1{,}25)\Rightarrow \heva{&u_m=0{,75}\\&u_{m+1}=1{,}25\\&n_m=32\\&n_{m-1}=25\\&n_{m+1}=14}$.\\
Do đó $M_0=0{,}75+\dfrac{32-25}{(32-25)+(32-14)}\cdot (1{,}25-0{,}75)=0{,}89$.
\itemch{\bf Sai}\\
Cỡ mẫu $n=90$ nên trung vị là $\dfrac{1}{2}(x_{45}+x_{46})\in [0{,}75;1{,}25)$.\\
Do đó $M_e=0{,}75+\dfrac{\dfrac{90}{2}-25}{32}\cdot 0{,}5=1{,}0625$.\\
Vậy $50\%$ số học sinh được đo có độ cận không vượt quá $1{,}0625$ độ.
\itemch{\bf Đúng}\\
Cỡ mẫu $n=90$ nên tứ phân vị thứ nhất là $x_{23}\in [0{,}25;0{,}75)$.\\
Do đó $Q_1=0{,}25+\dfrac{\dfrac{90}{4}-0}{25}\cdot 0{,}5=0{,}7$.\\
Tứ phân vị thứ ba là $x_{68}\in [1{,}25;1{,}75)$.\\
Do đó $Q_3=1{,}25+\dfrac{\dfrac{3\cdot90}{4}-(25+32)}{14}\cdot 0{,}5=1{,}625$.\\
Vậy $50\%$ số học sinh được đo có độ cận từ $0{,}7$ độ đến $1{,}625$ độ.
\end{itemchoice}
}
\end{ex}

\caukq
\begin{ex}%[1D1V5-3]
Giả sử một vật dao động điều hòa xung quanh vị trí cân bằng theo phương trình $x=2\cos\left(5t-\dfrac{\pi}{6} \right)$. Ở đây, thời gian $t$ tính bằng giây và quãng đường $x$ tính bằng centimét. Hãy cho biết trong khoảng thời gian từ $0$ đến $6$ giây, vật đi qua vị trí cân bằng bao nhiêu lần?
\shortans[6]{$9$}
\loigiai{
Vật đi qua vị trí cân bằng thì $x=0$
\begin{eqnarray*}
&\Leftrightarrow&2\cos\left(5t-\dfrac{\pi}{6} \right)=0\\
&\Rightarrow&5t-\dfrac{\pi}{6}=\dfrac{\pi}{2}+k\pi\\
&\Rightarrow&t=\dfrac{2\pi}{15}+k\dfrac{\pi}{5},\,k\in \mathbb{Z}.	
\end{eqnarray*}	
Trong khoảng thời gian từ $0$ đến $6$ giây 
\begin{eqnarray*}
&\Rightarrow&0\le t\le 6\\
&\Leftrightarrow&0\le\dfrac{2\pi}{15}+k\dfrac{\pi}{5}\le 6\\
&\Leftrightarrow &-\dfrac{2}{3}\le k\le 8{,}88\\
&\Rightarrow&k\in \left\lbrace 0;1;\ldots; 8 \right\rbrace.
\end{eqnarray*}	
Vậy trong khoảng thời gian từ $0$ đến $6$ giây, vật đi qua vị trí cân bằng $9$ lần.
}
\end{ex}

\begin{ex}%[1D2V2-7]
Một gia đình cần khoan một cái giếng để lấy nước. Họ thuê một đội khoan giếng đến để khoan giếng nước. Biết giá của một mét khoan đầu tiên là $75\,000$ nghìn đồng, kể từ mét khoan thứ hai giá mỗi mét khoan tăng lên $6\,000$ nghìn đồng so với giá của mét khoan trước đó. Biết cần phải khoang sâu $80\mathrm{\,m}$ thì mới có nước. Vậy gia đình cần phải trả bao nhiêu tiền để khoan cái giếng đó? (kết quả làm tròn đến hàng đơn vị của triệu đồng).
\shortans[6]{$25$}
\loigiai{
Giá tiền khoan mỗi mét giếng (bắt đầu từ mét đầu tiên) lập thành cấp số cộng $(u_n)$ có số hạng đầu $u_1=75\,000$ và công sai $d=6\,000$.\\
Do cần khoan $80\mathrm{\,m}$ nên tổng số tiền cần trả là\\ $$S_{80}=u_1+u_2+\ldots+u_{80}=\dfrac{80}{2}\left(2\cdot 75\,000+79\cdot 6\,000 \right)=24\,960\,000\approx 25\,\text{(triệu đồng)}.$$
}
\end{ex}

\begin{ex}%[1D2V3-8]
Bạn Lan thả một quả bóng cao su từ độ cao $12\mathrm{\,m}$ theo phương thẳng đứng. Mỗi khi chạm đất nó lại nảy lên theo phương thẳng đứng với độ cao bằng $\dfrac{2}{3}$ độ cao trước đó. Tính tổng quãng đường quả bóng đi được cho đến khi quả bóng dừng hẳn.
\shortans[6]{$60$}
\loigiai{
Các quãng đường khi bóng đi xuống tạo thành một cấp số nhân lùi vô hạn với số hạng đầu $u_1=12$ và công bội $q=\dfrac{2}{3}$.\\
Do đó tổng quảng đường mà quả bóng đi xuống là $S=\dfrac{u_1}{1-q}=\dfrac{12}{1-\dfrac{2}{3}}=36\mathrm{\,m}$.\\
Tổng quãng đường của quả bóng đi được cả lên và xuống đến khi quả bóng dừng hẳn là $$2S-12=2\cdot 36-12=60\mathrm{\,m}.$$
}
\end{ex}

\begin{ex}%[1D3V3-3]
Tìm giá trị nhỏ nhất của $a$ để hàm số $f(x)=\heva{&\dfrac{\sqrt{3x-2}-2}{x-2}&\text{khi}&\quad x> 2\\&a^2x+\dfrac{1}{4}&\text{khi}&\quad x\le 2}$ liên tục tại $x=2$.
\shortans[6]{$-0{,}5$}
\loigiai{
Ta có $\lim\limits_{x\to 2^{-}}f(x)=f(2)=2a^2+\dfrac{1}{4}$.\\
\begin{eqnarray*}
\lim\limits_{x\to 2^{+}}f(x)&=&\lim\limits_{x\to 2^{+}}\dfrac{\sqrt{3x-2}-2}{x-2}\\
&=&\lim\limits_{x\to 2^{+}}\dfrac{3(x-2)}{(x-2)\left(\sqrt{3x-2}+2\right)}\\
&=&\lim\limits_{x\to 2^{+}}\dfrac{3}{\sqrt{3x-2}+2}\\
&=&\dfrac{3}{4}.
\end{eqnarray*}	
Để hàm số liên tục tại $x=2\Leftrightarrow f(2)=\lim\limits_{x\to 2^{-}}f(x)=\lim\limits_{x\to 2^{+}}f(x)$.\\
Khi đó $2a^2+\dfrac{1}{4}=\dfrac{3}{4}\Leftrightarrow a=\pm 0{,}5$.\\
Vậy $a_{min}=-0{,}5$.
}
\end{ex}

\begin{ex}%[1H4V6-5]
Cho tứ diện $ABCD$. Gọi $G_1, G_2, G_3$ lần lượt là trọng tâm của các tam giác $ABC$, $ACD$ và $ADB$. Diện tích tam giác $BCD$ bằng $k$ lần diện tích thiết diện tạo bởi mặt phẳng $(G_1G_2G_3)$ với tứ diện $ABCD$. Tính giá trị của $k$.
\shortans[6]{$2{,}25$}
\loigiai{
\begin{center}
	\begin{tikzpicture}[scale=1,line join=round,font=\footnotesize,line cap=round,thick]
		\coordinate (A) at (0.5,5);
		\coordinate (B) at (-2,0);
		\coordinate (C) at (2,-2);
		\coordinate (D) at (4,0);
\path ($(B)!0.5!(C)$) coordinate (M);
\path ($(C)!0.5!(D)$) coordinate (N);
\path ($(B)!0.5!(D)$) coordinate (P);
\path ($(A)!2/3!(M)$) coordinate (G_1);
\path ($(A)!2/3!(N)$) coordinate (G_2);
\path ($(A)!2/3!(P)$) coordinate (G_3);
\path ($(A)!2/3!(B)$) coordinate (B');
\path ($(A)!2/3!(C)$) coordinate (C');
\path ($(A)!2/3!(D)$) coordinate (D');
		\draw(A)--(B) (A)--(C) (A)--(D) (C)--(D) (B)--(C)(A)--(M)(A)--(N) (B')--(C')--(D');
		\draw[dashed,thin](B)--(D)(A)--(P) (G_1)--(G_2)--(G_3)--cycle (M)--(N)--(P)--cycle (B')--(D');
		\foreach \i/\g in {A/90,B/180,C/-90,D/0,M/210,N/-45,P/-90,G_1/180,G_2/-15,G_3/140,B'/180,C'/-20,D'/45}{\draw[fill=black](\i) circle (1.5pt) ($(\i)+(\g:3mm)$) node[scale=1]{$\i$};}
	\end{tikzpicture}
\end{center}	
Gọi $M$, $N$, $P$ lần lượt là trung điểm của các cạnh $BC$, $CD$, $BD$.\\
Ta có $\dfrac{AG_1}{AM}=\dfrac{AG_2}{AN}=\dfrac{AG_3}{AP}=\dfrac{2}{3}\Rightarrow \heva{&G_1G_2\parallel MN\\&G_2G_3\parallel NP}\Rightarrow (G_1G_2G_3)\parallel (BCD)$.\\
Qua $G_1$ kẻ $B'C'\parallel BC$ ($B'\in AB$, $C'\in AC$).\\
Qua $G_2$ kẻ $C'D'\parallel CD$ ($D'\in AD$).\\
Khi đó thiết diện của hình chóp cắt bởi mặt phẳng $(G_1G_2G_3)$ là $\triangle B'C'D'$.\\
Ta có $\dfrac{B'C'}{BC}=\dfrac{AG_1}{AM}=\dfrac{2}{3}\Rightarrow \triangle B'C'D'\backsim \triangle BCD$ theo tỉ số $\dfrac{2}{3}$.\\
Do đó $\dfrac{S_{\triangle BCD}}{S_{\triangle B'C'D'}}=\left(\dfrac{3}{2}\right)^2=\dfrac{9}{4}=2{,25}$.
}
\end{ex}

\begin{ex}%[1D5V2-3]
Thời gian (đơn vị: phút) di chuyển đến trường của nhóm học sinh trường THPT A được tổng hợp trong bảng sau:
\begin{longtable}{|c|c|c|c|c|c|c|c|}
\hline
Thời gian (phút)&$[15;20)$&$[20;25)$&$[25;30)$&$[30;35)$&$[35;40)$&$[40;45)$&$[45;50)$\\
\hline
Số học sinh&$6$&$14$&$25$&$37$&$13$&$9$&$21$\\
\hline
\end{longtable}
Tìm khoảng tứ phân vị của mẫu số liệu ghép nhóm trên (kết quả làm tròn đến hàng phần mười).
\shortans[6]{$12{,}3$}
\loigiai{
Cỡ mẫu $n=125$.\\
Gọi $x_1,x_2,\ldots,x_{125}$ lần lượt là thời gian di chuyển đến trường của $125$ học sinh xếp theo thứ tự không giảm.\\
Tứ phân vị thứ nhất của mẫu số liệu $x_1,x_2,\ldots,x_{125}$ là $Q_1=\dfrac{1}{2}(x_{31}+x_{32})\in [25;30)$.\\
Do đó giá trị của tứ phân vị thứ nhất là $$Q_1=25+\dfrac{\dfrac{125}{4}-(6+14)}{25}\cdot(30-25)=27{,}25\,\text{(phút)}.$$
Tứ phân vị thứ ba của mẫu số liệu $x_1,x_2,\ldots,x_{125}$ là $Q_3=\dfrac{1}{2}(x_{94}+x_{95})\in [35;40)$.\\
Do đó giá trị của tứ phân vị thứ ba là $$Q_3=35+\dfrac{\dfrac{3\cdot 125}{4}-(6+14+25+37)}{13}\cdot(40-35)\approx 39{,}52\,\text{(phút)}.$$
Khoảng tứ phân vị của mẫu số liệu $\Delta Q=Q_3-Q_1=39{,}52-27{,}25\approx 12{,}3$ (phút).
}
\end{ex}